\appendix
\onecolumngrid

\section{Notation}
\label{app:notation}

Table~\ref{tab:symbols} defines every symbol of this paper. Symbols marked ``taken'' in the
provenance documents are those of the rest of PROJECT spacetimes, used here with exactly the meaning
that they have there (the meaning is copied from the project's master registries, with a note on the
use here where needed); all others are new. Every new symbol carries the suffix NG, except the elements
of the algebras (sans-serif type), the unit matrices, the standard number sets and operation signs; the
Dirac matrices and their relatives keep the tags 1671 ($7+1$) and 1644 ($4+4$) of the project. The
machine checks NG-S1 to NG-S6 (Appendix~\ref{app:checks}) verify that no symbol is registered twice,
that the taken symbols carry the meanings of the project's master registries, that no new symbol
coincides with a symbol of those registries, that every name in the data files is registered, and that
no letter with another meaning in the project is used in a formula of this paper without its suffix.
Units: $\hbar=1$, the speed of light 1, and the momentum cutoff $\Lambda_{\rm NG}=1$ in
Secs.~\ref{sec:superfluid} and~\ref{sec:renorm} unless a unit is shown.

{\small
\begin{longtable}{p{0.24\textwidth}p{0.54\textwidth}>{\raggedright\arraybackslash\footnotesize}p{0.16\textwidth}}
\caption{The symbols of this paper, by group, with the plain-text names under which the programs, the
figures and the data files write them.}\label{tab:symbols}\\
\toprule
symbol & meaning & plain-text names\\
\midrule
\endfirsthead
\toprule
symbol & meaning (continued) & plain-text names\\
\midrule
\endhead
\bottomrule
\endfoot
{{GENRAW:symbols_ng_texrows}}
\end{longtable}
}

\emph{Plain-text names.} The programs, the figures and the data files write formulas in plain text,
with the names of the last column of Table~\ref{tab:symbols}. A name followed by a tag in brackets in
a check report (\texttt{[7+1]}, \texttt{[4+4]}, \texttt{[O\_NG]}, \texttt{[Os\_NG]}, \texttt{[1671]},
\texttt{[1644]}) refers to that spacetime or algebra. The following names have a meaning that depends on
the place, which is stated here: \texttt{i}, \texttt{j}, \texttt{k}, \texttt{l}, \texttt{il},
\texttt{jl}, \texttt{kl} are the octonion units $\mathsf e_1,\dots,\mathsf e_7$ in checks O-14a, O-14c,
K-6, in Figs.~\ref{fig:tables} and~\ref{fig:manifolds} and in the header of the file
\texttt{ng\_octonion\_tables.csv} (and \texttt{1} is $\mathsf e_0$), while
elsewhere \texttt{i} is the imaginary unit and \texttt{j} the index of a space direction;
\texttt{x}, \texttt{y}, \texttt{z} are elements $\mathsf x$, $\mathsf y$, $\mathsf z$ of the algebras in
the checks O, T and K-5a and in the figures of the algebras, while \texttt{x0}, \dots, \texttt{x7} and
\texttt{xA} are always the coordinates; \texttt{M} and \texttt{K} in check K-6 are
$\mathcal M^{3,1}_{\rm NG}$ and $\mathcal K^{1,3}_{\rm NG}$, while \texttt{K} elsewhere (in
\texttt{gamma\^{}K}, \texttt{Z\_K}, \texttt{n4}) is an index set; \texttt{L} in check K-6 is
$L_{\rm NG}(\mathsf e_4)$; \texttt{a} in \texttt{L\_e\_A}, \texttt{e\_A}, \texttt{gamma\^{}A} is an index
of the eight axes and \texttt{a} in \texttt{Sigma\^{}(ab)}, \texttt{gamma\^{}a}, \texttt{Delta\_0a} a
fibre index; \texttt{C+g[0, 5]} in check K-4 is $C^+_{\rm NG}\gamma^0_{\rm NG}\gamma^5_{\rm NG}$;
\texttt{m4\^{}2} and \texttt{m8\^{}2} in check K-8 are the squares of the four-dimensional and the
eight-dimensional mass; \texttt{c\_N} in check P-5 is the factor $\tfrac12$ (one field) or 1 (doublet)
in front of the integral of Eq.~(\ref{eq:omega}); \texttt{Omega''/2} in check P-9 is half the second
derivative of $\Omega_{\rm NG}$ with respect to $\lvert\Delta_{\rm NG}\rvert$; \texttt{u} in check R-4 is
$\varpi_{\rm NG}$; \texttt{2 Sigma\^{}(ab)} is $2\Sigma^{ab}_{\rm NG}$; a superscript \texttt{T} is the
transpose and \texttt{dagger} the Hermitian conjugate; \texttt{8D} and \texttt{4D} abbreviate
eight-dimensional and four-dimensional, \texttt{BdG} and \texttt{BCS} Bogoliubov-de Gennes and
Bardeen-Cooper-Schrieffer. The groups and algebras \texttt{so(N)}, \texttt{so(8)},
\texttt{so(7,1)}, \texttt{so(4,4)}, \texttt{su(2)}, \texttt{g2}, \texttt{gl(8)}, \texttt{O(N)},
\texttt{SO(3)}, \texttt{U(1)}, \dots\ are written as in the text, and \texttt{Sym(3)} is the
permutation group of three objects.

\section{The machine checks}
\label{app:checks}

The program \texttt{nogo} (Appendix~\ref{app:data}) verifies {{GEN:nchecks_rust}} statements on which
the results rest; the script \texttt{check\_symbols.py} verifies the six statements NG-S1 to NG-S6 on
the symbols, and \texttt{check\_self\_contained.py} that no document of this revision sends the reader to
another document for a definition. Each check prints one line: its label, the statement, and, after two
dashes, the numbers that decide it. Table~\ref{tab:checks} reproduces these lines as the programs print
them, in plain text (Appendix~\ref{app:notation} explains the names).

{\footnotesize
\begin{longtable}{p{0.08\textwidth}p{0.80\textwidth}p{0.06\textwidth}}
\caption{All machine checks of this paper, as the programs print them.}\label{tab:checks}\\
\toprule
label & statement and decisive numbers & result\\
\midrule
\endfirsthead
\toprule
label & statement and decisive numbers (continued) & result\\
\midrule
\endhead
\bottomrule
\endfoot
{{GENRAW:checks_rows_tex}}
\end{longtable}
}

\section{Programs, data, and how to reproduce every number}
\label{app:data}

\subsection{Programs}

All computations are done by the Rust program \texttt{nogo} in the folder \texttt{code/nogo-rs} of
\texttt{Revision\_NoGo}; the figures are drawn by a Python script, and the documents are assembled by
Python scripts that insert every number from the data files, so that no computed number is typed by
hand. The program has twelve modules: \texttt{oct} (the Cayley-Dickson algebras, their products,
norms, associators, derivations and left multiplications, the octonionic Dirac matrices and the
intertwiners; command \texttt{octonions}), \texttt{triality} (the system~(\ref{eq:lt}), the triality maps,
their characteristic polynomials, the commutants; command \texttt{triality}), \texttt{sig} (the Dirac
matrices of the project on both spacetimes, the products $\gamma^K_{\rm NG}$, the Hamiltonian matrices and
the pair channels), \texttt{kaluza} (the combined manifolds, the pair classification, the $\mathrm G_2$
locking, the Kaluza-Klein spectra; command \texttt{kaluza}), \texttt{superalg} (the superalgebras, their
graded Jacobi identities and positivity; command \texttt{superalgebra}), \texttt{pairs4d} (the
Bogoliubov-de Gennes matrices, the effective potential, the gap equation, the amplitude mode, the
stabilizers; command \texttt{pairs4d}), \texttt{renorm} (power counting and the Kaluza-Klein bubble;
command \texttt{renorm}), \texttt{exact} (exact rational arithmetic with 128-bit integers, null spaces,
ranks, characteristic polynomials), \texttt{integrate} (adaptive Gauss-Legendre quadrature with 20 and
40 points), \texttt{rsym} (real symmetric matrices and the cyclic Jacobi method), \texttt{checks} (the
reports) and \texttt{main} (the commands). It uses the library crate \texttt{code/spacetimes-rs} of the
repository for the Dirac matrices $\gamma^A_{1671}$ and $\gamma^A_{1644}$, complex numbers, integer
matrices and their rank modulo a prime, complex Hermitian diagonalization, Gauss-Legendre nodes,
Brent's root finder, a parallel map and the writing of data files.

Numerical methods. The structural statements (checks O, T, K, S, R-1 to R-3) are decided in exact
integer or rational arithmetic; ranks of large integer matrices are computed modulo a large prime,
which can only underestimate a rank, so a rank found equal to its maximum is exact. The graded Jacobi
identities use structure constants that are integers, halves and quarters, which binary floating-point
arithmetic represents exactly. The momentum integrals of Sec.~\ref{sec:superfluid} use adaptive
Gauss-Legendre quadrature (the rules of 20 and 40 points are compared on each interval, which is halved
until they agree to the absolute tolerance $10^{-15}$ or the relative tolerance $10^{-14}$) on panels
whose end points cluster geometrically around the Fermi momentum; there $\varepsilon_{\rm NG}-\mu_{\rm NG}$
is computed as $(q^2-(q^{\rm F}_{\rm NG})^2)/(\varepsilon_{\rm NG}+\mu_{\rm NG})$, without cancellation. The
gap equation is solved by Brent's method in the logarithm of the gap; gaps below
$10^{-12}\Lambda_{\rm NG}$ are reported as zero.

\subsection{Commands}

From the folder \texttt{Revision\_NoGo} of the repository (a Rust toolchain for the edition 2021,
Python 3 with NumPy and Matplotlib, and for the documents a TeX distribution with REVTeX 4.2, pandoc and
XeLaTeX):
\begin{verbatim}
cargo run --release --manifest-path code/nogo-rs/Cargo.toml -- all
python3 code/scripts/make_figures.py
bash code/scripts/check_all.sh
\end{verbatim}
The first command writes every data file of the folder \texttt{data} (in about a second); the
commands \texttt{octonions}, \texttt{triality}, \texttt{kaluza}, \texttt{superalgebra}, \texttt{pairs4d}
and \texttt{renorm} in place of \texttt{all} run one part. The second draws the thirteen figures. The
third runs both, then the checks of the symbols and of self-containedness, writes the fragments of the
documents, and builds this paper and the provenance documents.

\subsection{Data files}

{{GENRAW:data_files_tex}}
Every data file states its content and units in its header lines (\texttt{\# content:}), and every
column name is either a descriptive word or one of the names listed in the last column of
Table~\ref{tab:symbols}.

\subsection{Checksums}

SHA-256 checksums of the data files, the figures and the programs from which this paper was built:
{\footnotesize
\begin{verbatim}
{{GENRAW:data_sha256}}
{{GENRAW:figure_sha256}}
{{GENRAW:code_sha256}}
\end{verbatim}
}

\twocolumngrid
